% =====================================================================
%  Streaming Bayesian Parameter Assimilation for OpenHydroTwin
%  Algorithm & Implementation Reference
%
%  This document describes the algorithm ACTUALLY IMPLEMENTED in
%  DTStreamingMCMC.{h,cpp} and DTAssimilation.cpp, including the exact
%  plateau, quorum, and convergence criteria, with their default
%  parameter values. Where the code stubs a designed-but-unimplemented
%  feature, that is called out explicitly in Section~\ref{sec:status}.
%
%  Build:  pdflatex streaming_mcmc_algorithm.tex   (run twice for refs)
% =====================================================================
\documentclass[11pt]{article}

\usepackage[margin=1in]{geometry}
\usepackage{amsmath,amssymb}
\usepackage{booktabs}
\usepackage{enumitem}
\usepackage{xcolor}
\usepackage{algorithm}
\usepackage{algpseudocode}
\usepackage[hidelinks]{hyperref}
\usepackage{titlesec}

\newcommand{\code}[1]{\texttt{\small #1}}
\newcommand{\param}[1]{\textsf{#1}}
\definecolor{stubgray}{gray}{0.35}
\newcommand{\stub}[1]{{\color{stubgray}\textit{#1}}}

\title{\bfseries Streaming Bayesian Parameter Assimilation\\ for OpenHydroTwin\\[4pt]
\large Algorithm and Implementation Reference}
\author{Digital Twin Runner --- \code{DTStreamingMCMC}}
\date{\today}

\begin{document}
\maketitle

\begin{abstract}
This document specifies the streaming Markov-chain Monte Carlo (MCMC)
parameter-assimilation algorithm used by the OpenHydroTwin digital-twin
runner. It documents the target distribution (a prior combined with a
recency-weighted likelihood), the multi-chain Metropolis--Hastings
sampler, and---in full detail---the two criteria that govern the method:
the \emph{per-chain plateau classifier} (when a chain has stopped
burning in) and the \emph{cycle convergence certification} (when a
cycle's pooled samples may be published as a certified posterior). All
formulas and default constants below reflect the code as implemented;
unimplemented (stubbed) design features are flagged in
Section~\ref{sec:status}.
\end{abstract}

\tableofcontents

% =====================================================================
\section{Notation}\label{sec:notation}
% =====================================================================

Symbols used throughout, in order of first appearance. Defaults for the
tunable constants are collected in Section~\ref{sec:defaults}.

\begin{center}
\begin{tabular}{@{}ll@{}}
\toprule
Symbol & Meaning\\
\midrule
$\psi=(\theta,\sigma)$ & full parameter vector: physical params $\theta$ and error std.\ devs.\ $\sigma$\\
$n_p$ & number of parameters (dimension of $\psi$)\\
$N_c$ & number of chains in the ensemble (\param{number\_of\_chains})\\
$\pi_t(\psi)$ & target posterior at window end $t_{\text{now}}$ (streaming)\\
$\log\pi_t$ & log-posterior $=$ log-prior $+$ log-likelihood (Sec.~\ref{sec:target})\\
$t_{\text{now}}$ & current calibration-window end (OHQ day-serial)\\
$t_0$ & earliest observation in the record (overall start)\\
$t_{\text{start}}$ & scored-window start $=\max(t_0,\ t_{\text{now}}-W_c)$\\
$W_c$ & rolling-window length (\param{calibration\_window\_days})\\
$\delta$ & observation age $t_{\text{now}}-t_i$\\
$w(\delta)$ & recency (forgetting) kernel weight, Eq.~\eqref{eq:kernel}\\
$\Delta_0,\tau,\alpha$ & kernel constants (plateau length, scale, decay exponent)\\
$\sigma$ & per-observation error standard deviation (part of $\psi$)\\
$x,X$ & current / proposed chain position\\
$s_i$ & per-parameter proposal scale (\code{pertcoeff}$_i$)\\
$\log J$ & Hastings log-correction for log-normal proposals\\
$\log A$ & MH log-acceptance ratio\\
$W_{\max},\,W_{\text{eff}}$ & max / effective plateau-classifier window (\param{plateauWindow})\\
$g$ & sweeps elapsed so far in the cycle\\
$n$ & trace points used in the trend test, $n=\min(|\text{trace}|,W_{\text{eff}})$\\
$y_k$ & $k$-th log-posterior trace value, $y_k\equiv\log\pi_t(\psi^{(k)})$\\
$k,\bar k$ & step index and its window mean $(n-1)/2$\\
$b$ & least-squares slope of $y$ vs.\ $k$\\
$s$ & residual scatter of the trend fit\\
$T$ & total trend across the window, $T=|b|(n-1)$\\
$R$ & decision statistic (trend-to-scatter ratio) $R=T/s$\\
$a$ & accepted steps within the window (motion guard)\\
$f_P$ & plateaued fraction $=$ \#Plateaued chains $/\,N_c$\\
$W_s$ & inter-cycle stability ring length (\param{stabilityWindow})\\
$\psi^{(r)}_i$ & parameter $i$'s point estimate at ring entry $r$\\
$N$ & pooled-sample count in a cycle\\
$\Sigma$ & joint sample covariance of the pool\\
$\gamma_k,\ \tau_{\text{int}}$ & lag-$k$ autocovariance and integrated autocorr.\ time (ESS)\\
$C$ & realization-reservoir capacity (\param{realizationCount})\\
\bottomrule
\end{tabular}
\end{center}

% =====================================================================
\section{Overview and terminology}
% =====================================================================

The runner alternates between two activities on a wall-clock schedule:

\begin{description}[leftmargin=1.4em]
  \item[Forward stepping.] The model is advanced through simulated time
  (\code{forward\_advance}) and a short forecast is produced
  (\code{forward\_forecast}), accumulating observations from the truth
  source into a rolling buffer.
  \item[Calibration cycle.] Every \param{poll\_interval} of simulated
  time (default \code{3day}), the accumulated observations trigger one
  \emph{calibration cycle} (\code{assim\_calibration}): a bounded run of
  the streaming sampler that re-estimates the posterior of the model
  parameters and emits a calibrated model snapshot for the forward loop
  to adopt.
\end{description}

A \emph{cycle} is one invocation of the sampler
(\code{DTStreamingMCMC::runCycle}). Within a cycle the sampler runs an
ensemble of $N_c$ Metropolis--Hastings \emph{chains} in parallel until a
wall-clock deadline. Each chain passes through two \emph{phases}:
\textbf{Climbing} (burn-in; its draws are disposable) and
\textbf{Plateaued} (stationary; its draws are retained). Only retained
draws are ever \emph{pooled} into the posterior estimate.

\paragraph{Governing design principle.} Selection \emph{among} chains is
free during the disposable Climbing phase and forbidden once a chain
Plateaus. All acceleration machinery (seeding, and---when
enabled---culling and dissolution) touches only pre-plateau chains; the
pooled posterior is gated on convergence alone.

Let $\psi=(\theta,\sigma)$ denote the full parameter vector of dimension
$n_p$, where $\theta$ are the physical model parameters and $\sigma$ the
per-observation error standard deviations (both are sampled jointly).

% =====================================================================
\section{The target distribution}\label{sec:target}
% =====================================================================

The sampler targets the log-posterior returned by the engine
(\code{posteriorLocal}), identical in value to the deployed GA mode's
objective:
\begin{equation}
  \log \pi_t(\psi) \;=\; \underbrace{\sum_{i=1}^{n_p}\log
  p_i(\psi_i)}_{\text{log-prior}}
  \;+\; \underbrace{\sum_{o}\log L_{o,t}(\psi)}_{\text{log-likelihood}}
  \;+\; \text{const},
\end{equation}
where the subscript $t$ marks the \emph{streaming} time dependence: the
likelihood re-weights observations by recency at the current window end
$t_{\text{now}}$.

\subsection{Recency (forgetting) kernel}
For an observation series with a ``Weighted Least Squared'' comparison
method and normal error structure, each observation of age
$\delta = t_{\text{now}} - t_i$ receives weight
\begin{equation}
  w(\delta) \;=\;
  \begin{cases}
    1, & \delta < \Delta_0,\\[4pt]
    \left(\,1 + \dfrac{\log(\delta/\Delta_0)}{\tau}\,\right)^{-\alpha},
      & \delta \ge \Delta_0,
  \end{cases}
  \label{eq:kernel}
\end{equation}
with per-observation kernel constants $\Delta_0$
(\param{kernel\_Delta0}, a full-weight plateau length), $\tau$
(\param{kernel\_tau}), and $\alpha$ (\param{kernel\_alpha}) taken from
the model file. Recent observations ($\delta<\Delta_0$) keep full
weight; older ones are down-weighted logarithmically. This is what makes
the target ``stream'': as $t_{\text{now}}$ advances, the effective data
window slides, and the posterior $\pi_t$ tracks a slowly moving target.

\subsection{Per-observation likelihood}
With kernel-weighted mean-squared error
$\mathrm{MSE}_w = \big(\sum_i w_i (y_i-\hat y_i)^2\big)/\sum_i w_i$,
effective count $W=\sum_i w_i$, and error standard deviation $\sigma$
(\param{error\_standard\_deviation}, itself calibratable), the engine
returns as objective the negative log-likelihood contribution
\begin{equation}
  -\log L_{o,t} \;=\; W\!\left(\frac{\mathrm{MSE}_w}{2\sigma^2}
  + \log\sigma\right),
  \label{eq:nll}
\end{equation}
so that $\sigma$ remains statistically identifiable when sampled jointly
with $\theta$. (For log-normal error structure the same form is applied
to log-transformed series.)

Because $\mathrm{MSE}_w$ is a weighted \emph{mean}, $\sum_i w_i r_i^2 = W\,
\mathrm{MSE}_w$, and Eq.~\eqref{eq:nll} is the Gaussian negative
log-likelihood $W\log\sigma + \sum_i w_i r_i^2/(2\sigma^2)$ with the
additive constant $\tfrac{W}{2}\log 2\pi$ dropped. The factor $2$ is what
makes $\hat\sigma$ the RMS residual: minimising Eq.~\eqref{eq:nll} over
$\sigma$ gives $\hat\sigma^2 = \mathrm{MSE}_w$. \emph{Historical note:} the
engine omitted this factor until 2026-07, returning
$W(\mathrm{MSE}_w/\sigma^2 + \log\sigma)$, whose minimiser is
$\hat\sigma = \sqrt2\,\mathrm{RMS}$ --- so every calibrated
\param{error\_standard\_deviation} in runs predating the fix is inflated by
exactly $\sqrt2$. Profiling $\sigma$ out gives
$\text{const} + \tfrac{W}{2}\log \mathrm{MSE}_w(\theta)$ under \emph{both}
forms, so the correction shifts $\sigma$ only and leaves inference on
$\theta$ unchanged.

The sampler stores honest log-posteriors:
the log-prior and this likelihood only; the Hastings/Jacobian correction
of Section~\ref{sec:proposal} is \emph{not} folded in.

\subsection{Rolling calibration window and initial
condition}\label{sec:window}
The likelihood above runs over a \emph{bounded} data window, not the
entire record. Let $t_0$ be the earliest observation and $t_{\text{now}}$
the latest. With window length $W_c$ (\param{calibration\_window\_days}),
the scored window is
\begin{equation}
  [t_{\text{start}},\,t_{\text{now}}],\qquad
  t_{\text{start}} = \max\!\big(t_0,\; t_{\text{now}} - W_c\big).
\end{equation}
Observations older than $t_{\text{start}}$ are \emph{excluded from the
solve entirely} (not merely down-weighted): the model is integrated only
over $[t_{\text{start}}, t_{\text{now}}]$. This is what keeps the
per-cycle cost bounded --- without it, each cycle re-integrates the whole
growing record and the number of sweeps affordable under a fixed budget
collapses as the run lengthens. The recency kernel of
Eq.~\eqref{eq:kernel} then applies \emph{within} the window. Setting
$W_c \le 0$ recovers the full-record window ($t_{\text{start}}=t_0$).

\paragraph{Initial condition via a MAP spin-up.} When
$t_{\text{start}} > t_0$ the solve no longer starts from the known cold
(script) initial state, so the model state at $t_{\text{start}}$ must be
established. Each cycle performs \emph{one} forward solve from $t_0$ to
$t_{\text{start}}$ using the latest maximum-a-posteriori (MAP) parameter
set, started from the cold initial state, and captures its end state.
That state seeds \emph{every} chain, so all chains share a single,
MAP-consistent initial condition at the window start; each chain then
integrates $[t_{\text{start}}, t_{\text{now}}]$ from it under its own
proposed parameters. The spin-up is one solve per cycle --- negligible
against the $N_c \times$(sweeps) solves of the sampling loop --- and its
cost is dominated by that loop even though it grows with $t_{\text{start}}$.
A spin-up failure degrades safely to the full-record window.

% =====================================================================
\section{The Metropolis--Hastings ensemble}
% =====================================================================

\subsection{Chains and parallelism}
The ensemble holds $N_c=$ \param{number\_of\_chains} chains (from the
model's \code{MCMC} settings). Each cycle sweeps all chains in parallel
(OpenMP, \param{numberOfThreads} threads); within a sweep, chain $c$
performs one MH step followed by one plateau classification, touching
only its own state, so the sweep is race-free. Each chain owns a
private, pristine copy of the \code{System}, so no two threads ever read
the same model object concurrently. Per-chain RNGs are seeded serially
from a master seed, making a run reproducible for a fixed seed regardless
of thread scheduling.

\subsection{Proposal mechanism}\label{sec:proposal}
For chain position $x$, the proposal $X$ is drawn coordinate-wise
(\code{proposeFrom}) according to each parameter's prior type, using
per-parameter scale $s_i=$ \code{pertcoeff}$_i$ and $Z\sim\mathcal N(0,1)$:
\begin{align}
  \text{additive (normal/uniform prior):}\quad & X_i = x_i + s_i Z,\\
  \text{multiplicative (log-normal prior):}\quad & X_i = x_i \,e^{\,s_i Z}.
\end{align}
The multiplicative (log-normal random-walk) proposal is asymmetric; its
Hastings log-correction is accumulated explicitly,
\begin{equation}
  \log J \;=\; \sum_{i \in \text{log-normal}} \big(\log X_i - \log x_i\big),
\end{equation}
and is $0$ for an all-additive parameter set.

\paragraph{Adaptive-covariance mode.} When
\param{mcmc\_proposal\_mode}${}=$\code{"covariance"}, the coordinate-wise
walk above is replaced by a proposal preconditioned by the accumulated
posterior covariance $\hat\Sigma = LL^{\!\top}$ (Cholesky $L$, ridge
$\varepsilon I_d$), $X = x + \kappa L z$ with $z\sim\mathcal N(0,I_d)$ and
$\kappa_0 = 2.38/\sqrt d$. The factorisation is carried out in
\emph{proposal space} --- log coordinates for log-normal priors --- so the
multiplicative form and its Hastings correction are unchanged. Until
$\hat\Sigma$ has accumulated weight $\ge 2d$ the proposal falls back to the
coordinate-wise walk. See \code{proposal\_adaptation\_methods.tex} for the
derivation and the accumulation rule.

\subsection{Acceptance rule}
The proposal is accepted (\code{stepChain}) by the standard MH rule in
log space, with the Hastings correction entering the ratio (never the
stored $\log\pi$). Non-finite proposals always reject. With
$\log A = \big(\log\pi_t(X) + \log J\big) - \log\pi_t(x)$:
\begin{equation}
  \text{accept if } \log A \ge 0,\quad
  \text{else accept with probability } e^{\log A}
  \ \ (\text{i.e. } \log U < \log A,\ U\sim\mathcal U(0,1)).
\end{equation}
A step costs exactly one forward solve---the dominant cost of the method.

% =====================================================================
\section{Plateau classifier (burn-in detection)}\label{sec:plateau}
% =====================================================================

The classifier (\code{classifyChain}) decides, from a chain's own recent
\emph{trajectory} of log-posterior values, whether it is still climbing
toward the typical set or has become stationary. It is
\textbf{trajectory-based, never level-based}: only the trend of a chain's
$\log\pi$ trace matters, never its absolute height---so a chain that is
flat-but-moving on a low-density ridge or secondary mode still qualifies
as plateaued.

\subsection{Effective window}
Let $W_{\max}=$ \param{plateauWindow} (default $60$) be the maximum trace
length, and let $g$ be the number of sweeps elapsed so far. The effective
window is
\begin{equation}
  W_{\text{eff}} \;=\; \min\!\big(W_{\max},\ \max(10,\ \lfloor g/2\rfloor)\big),
\end{equation}
so that late, short cycles (whose sweep budget shrinks as the data record
grows) can still fill a window and classify instead of being stuck
``climbing'' for lack of data. The trend is tested over the last
$n=\min(|\text{trace}|, W_{\text{eff}})$ points.

\subsection{Guards before classification}
A chain is forced to \textbf{Climbing} if either:
\begin{enumerate}[nosep,leftmargin=1.6em]
  \item $n < \min(\param{minStepsBeforeClassify},\,W_{\text{eff}})$
  (default floor $20$): too little data; or
  \item any value in the window is non-finite or a sentinel
  ($\le -10^{299}$): the chain is parked in an infeasible region.
\end{enumerate}

\subsection{Trend and scatter}
Each chain records a trace of its \emph{log-posterior} values: one entry
$y_k \equiv \log\pi_t\!\big(\psi^{(k)}\big)$ is appended per Metropolis
step ($\code{chain.trace.push\_back(chain.logp)}$), holding the newly
accepted value on an accept and the retained current value on a reject.
Over the last $n$ trace points $y_0,\dots,y_{n-1}$ (with the step index
$k$ as the regressor, $\bar k = (n-1)/2$), fit a least-squares line to
$y$ versus $k$ and measure its residual scatter:
\begin{equation}
  b = \frac{\sum_{k}(k-\bar k)(y_k-\bar y)}{\sum_k (k-\bar k)^2},
  \qquad
  s = \sqrt{\frac{1}{\max(n-2,\,1)}\sum_k\big[(y_k-\bar y)-b(k-\bar k)\big]^2}.
\end{equation}
The \emph{total trend across the window} is $T = |b|\,(n-1)$, i.e.\ the
fitted rise/fall in $\log\pi$ over the whole window, and the decision
statistic is the trend-to-scatter ratio $R = T/s$.

\subsection{Motion (stagnation) guard}
A window carried by too few \emph{accepted} moves is a \emph{stuck} chain,
not a converged one: its trace is flat because nothing moved. Let $a$ be
the number of accepted steps within the window. The chain is
\emph{moving} iff
\begin{equation}
  a \;\ge\; \max\!\big(1,\ \lfloor \param{minAcceptedFraction}\cdot n\rfloor\big),
  \qquad \param{minAcceptedFraction}=0.05.
\end{equation}
If not moving, the chain is \emph{not} plateaued (it is dissolution's
territory---see Section~\ref{sec:status}, where dissolution is noted as
currently stubbed).

\subsection{Decision with hysteresis}
A moving chain is classified using a two-level (hysteretic) threshold on
$R$ to prevent flapping when $R$ hovers near the boundary. With
$\param{plateauSlopeThreshold}=1.0$ and
$\param{plateauRevertFactor}=3.0$:
\begin{equation}
  \text{Plateaued} \iff
  \text{moving} \ \wedge\
  \begin{cases}
    R < \param{plateauSlopeThreshold}, & \text{if currently Climbing (declare)},\\[2pt]
    R < \param{plateauSlopeThreshold}\cdot\param{plateauRevertFactor}, & \text{if currently Plateaued (revert only above)}.
  \end{cases}
  \label{eq:plateau}
\end{equation}
In words: \emph{a chain plateaus when the trend it has drifted over the
window is less than one within-window standard deviation of noise}, and
once plateaued it reverts to Climbing only if the trend grows past three
such standard deviations. A ``Plateaued $\to$ Climbing'' reversion is
logged as ``target moved under the chain''; already-retained samples are
kept (they were drawn while genuinely stationary), but retention pauses
until the chain re-plateaus. The special case $s=0$ with accepted moves
(effectively impossible for a continuous target) is classified
conservatively as not plateaued.

% =====================================================================
\section{Pooling and the plateaued fraction}
% =====================================================================

A sample is \emph{retained} (appended to its chain's sample store) on a
step \emph{only while that chain is classified Plateaued}
(\code{stepChain}); Climbing samples are never stored. At cycle end,
\code{assemblePool} concatenates every chain's retained samples with
\textbf{equal weight}---the posterior shape is carried by sample density
alone, and burn-in exclusion has already happened at retention time. Only
the current cycle's samples are ever pooled; distinct per-cycle targets
are never mixed.

The \emph{plateaued fraction} is the ensemble-wide readiness statistic
\begin{equation}
  f_P \;=\; \frac{\#\{c : \text{phase}(c)=\text{Plateaued}\}}{N_c}.
\end{equation}

% =====================================================================
\section{Convergence certification}\label{sec:convergence}
% =====================================================================

At cycle end (\code{runCycle}) the result is certified \textbf{FULL}
(publishable posterior) or \textbf{PROVISIONAL} (point product only). The
\emph{primary} criterion is a within-cycle quorum; an \emph{optional},
weaker inter-cycle stability path may additionally certify:
\begin{equation}
  \texttt{converged} \;=\; \text{(byQuorum)} \ \lor\ \text{(byStability)}.
\end{equation}
All thresholds below are configurable (\param{mcmc\_quorum\_fraction},
\param{mcmc\_stability\_*}); the values shown are code defaults.

\subsection{Path 1 --- within-cycle quorum (primary)}
\begin{equation}
  \text{byQuorum} \iff f_P \ge \param{quorumFraction},
  \qquad \param{quorumFraction}=0.5\ \text{(default)}.
\end{equation}
At least a fraction $q$ of the chains have \emph{independently} reached
stationarity within this single cycle. This is a purely within-cycle test,
so it is \textbf{compatible with genuine parameter drift}: a certified
cycle whose estimate differs from the previous one records a real, tracked
change, not a convergence failure. Lowering $q$ certifies a larger
proportion of cycles (fewer plateaued chains required); it is the natural
knob for trading certification rate against per-cycle evidence.

\subsection{Path 2 --- inter-cycle stability (optional fallback)}
This weaker path targets the regime where a within-cycle quorum cannot
form (late cycles with few sweeps) but the point estimate has
\emph{settled across cycles}. A ring buffer holds the last
$W_s=\param{stabilityWindow}$ (default $5$) cycle point estimates, and the
stream is \emph{stable} (\code{streamStable}) iff the ring is full and
\emph{every} parameter's spread across the ring is small relative to its
recent mean:
\begin{equation}
  \text{streamStable} \iff
  \max_{i}\ \frac{\max_r \psi_i^{(r)} - \min_r \psi_i^{(r)}}
  {\max(|\bar\psi_i|,\,10^{-9})} \;\le\; \param{stabilityTol},
  \qquad \param{stabilityTol}=0.02\ \text{(default)},
\end{equation}
with the certified condition
\begin{equation}
  \text{byStability} \iff \param{stabilityEnabled}\ \wedge\
  (\neg\,\text{byQuorum}) \ \wedge\ \text{streamStable} \ \wedge\
  f_P \ge \param{stabilityMinPlateaued}\ (=0.1).
\end{equation}

\paragraph{Caveat --- stationarity vs.\ drift.} This path
\emph{presumes an approximately stationary target}: it certifies only when
the point estimate stops moving. That is at odds with \emph{detecting
parameter drift}, a core aim of the streaming method --- a genuinely
drifting parameter would be (correctly) tracked cycle to cycle, yet this
test would read the motion as non-convergence and withhold certification.
It is therefore \textbf{off by default in drift studies}
(\param{stabilityEnabled}=false, config key
\param{mcmc\_stability\_enabled}), leaving the within-cycle quorum as the
sole criterion. It is retained as an opt-in fallback for
stationary-target deployments. A further subtlety: the point estimate is
the sampled MAP (the mode), which for a \emph{non-identifiable} parameter
wanders randomly inside its (stable, wide) posterior, so the mode-based
spread overstates drift for such parameters --- another reason the path is
unsuitable when non-identifiability is expected.

Either way, certification relaxes only \emph{how many} plateaued chains
are required to trust the pool---never \emph{which} samples pool. Only the
current cycle's plateaued samples are ever used; no cross-cycle mixing
occurs. The certification source (\code{"quorum"}, \code{"stability"}, or
\code{""}) is recorded in the snapshot and history for auditing.

\subsection{Provisional cycles}
If neither path holds, the cycle is \textbf{provisional}: it still
publishes a point estimate and (as implemented) the dispersion of its
partial pool, but that dispersion is an \emph{uncertified} estimate, not
a certified posterior width. Provisional cycles are the expected regime
while the ensemble migrates toward a relocated target; each is expected
to leave the ensemble better located than it found it (see
Section~\ref{sec:seeding}).

% =====================================================================
\section{Proposal-scale adaptation}
% =====================================================================

Every \param{adaptationBlock} (default $20$) sweeps, \code{adaptProposalScale}
measures the pooled acceptance rate $\rho$ over the just-closed block and
nudges every per-parameter scale toward the configured target
\param{acceptance\_rate} using the base engine's multiplicative rule with
factor $\kappa=\param{purt\_change\_scale}$:
\begin{equation}
  s_i \leftarrow
  \begin{cases}
    s_i/\kappa, & \rho > \param{acceptance\_rate}\ \ (\text{accepting too often}\Rightarrow\text{enlarge steps}),\\
    s_i\cdot\kappa, & \rho \le \param{acceptance\_rate}\ \ (\text{shrink steps}).
  \end{cases}
\end{equation}
\textbf{Freeze on quorum.} Once a quorum holds, adaptation is
\emph{frozen}: continued kernel changes while samples are being retained
would break the invariance the pooled posterior rests on. Acceptance
bookkeeping continues so the cycle acceptance rate stays accurate, but
$s_i$ is held fixed---post-plateau draws are taken under a fixed kernel.

\textbf{In adaptive-covariance mode} the same block rule drives the single
global scalar $\kappa$ instead of the vector $s$, clamped to
$[10^{-4},10^{3}]$; the proposal \emph{shape} is then set by $\hat\Sigma$
and only its overall size is adapted, which is what stops one scalar from
having to serve a sharp and a broad direction at once. In this mode $s_i$
is left at its initialisation value and is not a meaningful diagnostic.

\paragraph{Two defects fixed in 2026-07 (both silent).} They are recorded
here because a run predating the fix will show the pathology and no error.
\begin{enumerate}[leftmargin=1.6em,nosep]
  \item \textbf{$\hat\Sigma$ did not survive the cycle.} \code{DTStreamingMCMC}
  is constructed fresh per calibration cycle, so the accumulator --- folded
  in \emph{after} the sampling loop and never persisted --- was rebuilt from
  one cycle and discarded. The readiness test $W\ge 2d$ was therefore
  satisfied only moments before destruction, and the covariance proposal ran
  for exactly zero proposals in a 244-cycle run. $\hat\Sigma$, $W$ and
  $\kappa$ now round-trip through \code{posterior\_latest.json} and are
  restored (and refactorised) in \code{initializeCycle}, so the kernel is
  armed from the first sweep.
  \item \textbf{$\kappa$ was pinned at its floor.} $\kappa$ initialised to
  $0$ and its lazy assignment lived in \code{refactorProposalCholesky},
  reached only at cycle end; every adaptation block computed
  $0\cdot\param{purt\_change\_scale}=0$ and the clamp lifted it to
  $10^{-4}$ --- about $8000\times$ below $\kappa_0$. It is now initialised
  in \code{initializeCycle}, before any adaptation block can run.
\end{enumerate}
Together these meant \code{"covariance"} silently degraded to the legacy
walk \emph{with adaptation disabled}: $s_i$ frozen at initialisation for the
whole run, while $\kappa$ was tuned diligently and read by nothing.

% =====================================================================
\section{Cross-cycle seeding}\label{sec:seeding}
% =====================================================================

Each cycle seeds its chains from the previous cycle's product
(\code{chooseSeedMode}), banking cross-cycle progress:
\begin{center}
\begin{tabular}{@{}lll@{}}
\toprule
Condition & Seed mode & Status\\
\midrule
no previous snapshot        & \code{ColdStart}          & interim (see \S\ref{sec:status})\\
previous cycle provisional  & \code{ProvisionalResume}  & implemented\\
drift detected              & \code{RatioReseed}        & \stub{not implemented}\\
otherwise (stationary)      & \code{WarmStart}          & implemented\\
\bottomrule
\end{tabular}
\end{center}

\begin{description}[leftmargin=1.4em,nosep]
  \item[WarmStart.] Uniform, \emph{unweighted} draws with replacement
  from the previous cycle's pooled posterior. Weighting them by posterior
  value would apply the density twice and under-disperse the ensemble---a
  bias that compounds because each cycle seeds the next. Seed
  log-posteriors are re-evaluated against the \emph{current} target, not
  copied.
  \item[ProvisionalResume.] Each chain picks up exactly where it left off,
  so migration toward a relocated target accumulates across cycles rather
  than restarting. If the chain count changed (config edit), chains resume
  by uniform draw from the carried states.
  \item[ColdStart / RatioReseed.] See Section~\ref{sec:status}.
\end{description}

% =====================================================================
\section{Point estimates and posterior summaries}
% =====================================================================

\subsection{Point estimate for the forward loop}
The forward loop adopts the cycle's point estimate, taken as the
\emph{sampled MAP} (\code{posteriorMAP}): the highest-$\log\pi$ retained
sample when FULL, or the highest-posterior carried chain state when
provisional. The posterior mean (\code{posteriorMean}) is also available.
The point estimate is pushed onto the stability ring each cycle.

\subsection{Marginal and joint summaries}
For a pool of $N$ samples (\code{computeSummary}) the runner computes, per
parameter $i$: the mean $\bar\psi_i$; the joint sample covariance
$\Sigma$ (denominator $N-1$), whose diagonal gives the standard
deviations and whose off-diagonals expose parameter compensation
structure; and marginal percentiles by sorting each column and linearly
interpolating at rank $q(N-1)$ for
$q\in\{0.025,\,0.10,\,0.50,\,0.90,\,0.975\}$. The
$2.5/50/97.5$ percentiles form the $95\%$ credible interval reported to
the viewer.

\subsection{Effective sample size}
ESS (\code{effectiveSampleSize}) is computed \emph{per chain} and summed,
never on the flat concatenated pool (which would show false correlation
at chain seams). For a scalar chain of length $n\ge 8$ with autocovariance
$\gamma_k$, Geyer's initial monotone-sequence estimator pairs adjacent
lags $\Gamma_m=\gamma_{2m}+\gamma_{2m+1}$, sums them until the first
non-positive pair (with a monotone-decreasing clamp), and forms
\begin{equation}
  \tau_{\text{int}} = \frac{2\sum_m \Gamma_m - \gamma_0}{\gamma_0},
  \qquad
  \mathrm{ESS}_{\text{chain}} = \frac{n}{\tau_{\text{int}}}.
\end{equation}
Per-parameter ESS is summed across chains; the reported cycle ESS is the
\emph{minimum} across parameters---the worst-mixing direction bounds the
pool's reliability. Chains shorter than $8$ samples are treated as
independent (ESS $=n$).

% =====================================================================
\section{Posterior-predictive products}
% =====================================================================

\subsection{Realization credible band (reservoir sampling)}
The predictive band over the observed quantities is built \emph{without
any extra solves}. During sampling, whenever a plateaued chain accepts a
step, the modeled output series from that already-computed solve are
offered to a fixed-capacity reservoir (\code{offerToReservoir}) of size
\param{realizationCount} (default $100$). Reservoir sampling keeps a
uniform random subset of all offered plateaued realizations: the first
$C$ fill it, and the $m$-th offer thereafter replaces a random slot with
probability $C/m$. At cycle end (\code{produceRealizationCI}), for each
\emph{calibration} observation the $2.5/50/97.5$ percentile bracket across
the retained realizations is written. All calibration observations are
written to a single combined file with columns labelled
\code{<observation> | <percentile>}.

\subsection{Parameter posterior histograms}
\code{writePosteriorDistribution} writes, per parameter, a histogram over
the pooled samples with bin count $\mathrm{clamp}(\lfloor\sqrt N\rceil,
10, 60)$, plus a mean/stdev/$2.5/50/97.5$ CI table.

\subsection{Cumulative posterior samples}
For building smooth publication distributions, \code{appendPooledSamples}
appends every cycle's raw pooled draws to a single running file
(\code{posterior\_samples.csv}), one row per draw, tagged with its origin
\code{cycle}, \code{t\_now}, and \code{converged} flag. The per-cycle pool
is otherwise discarded after the summary/histogram are computed, so this
file is the only cumulative record of the individual draws. Because each
cycle targets its own (sliding-window) posterior, downstream code should
pool over a chosen set of cycles rather than the entire file; the tags
make that selection explicit. This is the mechanism that makes a usable
posterior out of small per-cycle pools: a single cycle retains only tens of
draws (of which a substantial fraction are the repeats a rejected step
re-pushes), so publication densities are built by pooling a window of
consecutive cycles within one regime.

\subsection{Post-hoc analysis dumps}
Two records exist only for offline analysis; neither is read back by the
sampler, and both are written on the \param{mcmc\_realization\_interval}
gate alongside the products above.

\code{writeProposalCovariance} emits
\code{proposal\_cov\_cycle\_<t>.txt}: the accumulated $\hat\Sigma$ in
proposal space, together with the marginal standard deviations and the
derived correlation matrix, plus $W$, $\kappa$ and the readiness flag in
the header. Without it the running estimate is visible only in the latest
snapshot, which each cycle overwrites, so its convergence over a run is
unrecoverable.

\code{writeChainTraces} emits \code{chain\_traces\_cycle\_<t>.csv} with
columns \code{chain, sweep, logp, accepted, phase}. \code{DTChainState::trace}
is a rolling deque capped at \param{plateauWindow}, so the climb-to-plateau
trajectory that the classifier of Section~\ref{sec:plateau} acts on is
otherwise destroyed as the cycle proceeds; this is its only record. The
recorded \code{phase} is the one in force when the step was taken, so a
\code{'P'} row is a step whose sample entered the pool.

Per-cycle scalars --- including $\kappa$, the covariance readiness flag and
weight, the seed mode, the distinct-draw count and the across-chain
$\log\pi$ spread --- are written to \code{posterior\_history.jsonl}
(Section~\ref{sec:status}).

% =====================================================================
\section{Cycle driver and wall-clock budget}
% =====================================================================

A cycle stops at \emph{whichever comes first}: a wall-clock deadline or a
sweep cap (\code{runCycle}). The deadline is derived from the calibration
cadence: budget $=$ \param{poll\_interval} wall-clock duration minus a
safety margin (or an explicit \param{mcmc\_budget}). The cap is
\param{mcmc\_max\_sweeps} (default $300$; $\le 0$ disables it, leaving the
deadline as the sole limit). With the bounded rolling window
(Section~\ref{sec:window}) keeping solve cost roughly constant, the sweep
cap is typically the effective control on cheap cycles and the deadline on
expensive ones. The deadline is checked between sweeps (every
\param{stepsPerClockCheck} sweeps); an in-flight sweep is never aborted, so
a cycle may overrun by at most one forward-solve duration, which the margin
absorbs. If the deadline is already past on entry, the cycle degrades to a
zero-sweep provisional result rather than failing.

\begin{algorithm}[h]
\caption{One calibration cycle (\code{runCycle}), as implemented}
\begin{algorithmic}[1]
\State if window start $> t_0$: MAP spin-up $t_0\!\to\!t_{\text{start}}$ $\Rightarrow$ shared initial condition (Section~\ref{sec:window})
\State choose seed mode from previous snapshot; seed and evaluate all $N_c$ chains
\While{now $<$ deadline \textbf{and} sweeps $<$ \param{maxSweeps}}
  \ForAll{chains $c$ in parallel}
     \State one MH step (one forward solve); classify $c$ (Section~\ref{sec:plateau})
     \State if $c$ plateaued and step accepted: retain sample; offer outputs to reservoir
  \EndFor
  \State \stub{(serial interventions: dissolution / cull --- currently stubs)}
  \If{sweep index $\bmod$ \param{adaptationBlock} $= 0$} adapt proposal scales (frozen after quorum) \EndIf
\EndWhile
\State certify: \texttt{converged} $\gets$ byQuorum $\lor$ byStability\ (stability optional; Section~\ref{sec:convergence})
\State assemble pool; compute summaries + ESS if pool non-empty
\State \Return cycle result (FULL or PROVISIONAL)
\end{algorithmic}
\end{algorithm}

% =====================================================================
\section{Default parameters}\label{sec:defaults}
% =====================================================================

Streaming-specific knobs live in \code{DTStreamingSettings}; the base MCMC
knobs (chains, threads, target acceptance, scale-change factor) are parsed
from the model file's \code{MCMC} settings object. Entries marked
\textsuperscript{$\dagger$} are set per deployment from \code{config.json}
(keys in parentheses); the rest are code defaults.

\begin{center}
\begin{tabular}{@{}lll@{}}
\toprule
Symbol / setting & Default & Role\\
\midrule
\multicolumn{3}{@{}l}{\textit{Rolling window \& driver (\S\ref{sec:window}, \S\ref{sec:convergence})}}\\
\param{calibrationWindowDays}\textsuperscript{$\dagger$} & $365$ & scored-window length (\param{calibration\_window\_days}); $\le0$ = full record\\
\param{maxSweeps}\textsuperscript{$\dagger$} & $300$ & per-cycle sweep cap (\param{mcmc\_max\_sweeps}); $\le0$ = deadline only\\
\midrule
\multicolumn{3}{@{}l}{\textit{Plateau classifier (\S\ref{sec:plateau})}}\\
\param{plateauWindow} & $60$ & max trace length for the trend test\\
\param{plateauSlopeThreshold} & $1.0$ & declare plateau when $R<$ this\\
\param{plateauRevertFactor} & $3.0$ & revert only when $R>$ threshold$\times$this\\
\param{minStepsBeforeClassify} & $20$ & below this a chain is always Climbing\\
\param{minAcceptedFraction} & $0.05$ & motion (anti-stagnation) floor\\
\midrule
\multicolumn{3}{@{}l}{\textit{Convergence (\S\ref{sec:convergence})}}\\
\param{quorumFraction}\textsuperscript{$\dagger$} & $0.5$ & plateaued fraction for byQuorum (\param{mcmc\_quorum\_fraction})\\
\param{stabilityEnabled}\textsuperscript{$\dagger$} & \code{true} & enable inter-cycle path (\param{mcmc\_stability\_enabled})\\
\param{stabilityWindow}\textsuperscript{$\dagger$} & $5$ & cycles compared for byStability\\
\param{stabilityTol}\textsuperscript{$\dagger$} & $0.02$ & max relative point-estimate drift\\
\param{stabilityMinPlateaued}\textsuperscript{$\dagger$} & $0.1$ & plateaued floor to trust a partial pool\\
\midrule
\multicolumn{3}{@{}l}{\textit{Driver \& adaptation}}\\
\param{adaptationBlock} & $20$ & sweeps between scale adaptations\\
\param{stepsPerClockCheck} & $1$ & sweeps between deadline checks\\
\midrule
\multicolumn{3}{@{}l}{\textit{Predictive products}}\\
\param{realizationCount} & $100$ & reservoir size (predictive band)\\
\param{realizationInterval}\textsuperscript{$\dagger$} & $10$ & publish band/dist every $N$th cycle (\param{mcmc\_realization\_interval})\\
\midrule
\multicolumn{3}{@{}l}{\textit{Designed but stubbed (\S\ref{sec:status})}}\\
\param{cullEnabled} & \code{true} & (cull check is a stub)\\
\param{cullMarginLogp} & $20.0$ & log-posterior margin for culling\\
\param{stuckThreshold} & $200$ & consecutive rejections for dissolution\\
\param{minAcceptanceRate} & $10^{-3}$ & dissolution acceptance floor\\
\param{coldStartMultiplier} & $20$ & $K=$ this $\times N_c$ prior candidates\\
\param{explorerFraction} & $0.03$ & chains seeded raw from the prior\\
\bottomrule
\end{tabular}
\end{center}

% =====================================================================
\section{Implementation status (what is and is not active)}\label{sec:status}
% =====================================================================

The header documents a fuller design (referred to as ``spec v2''); several
acceleration features are present as interfaces but currently stubbed. For
an accurate reading of runtime behaviour:

\begin{itemize}[leftmargin=1.4em]
  \item \textbf{Culling} (\code{cullCheck}), \textbf{dissolution}
  (\code{dissolutionCheck}), \code{cloneChain}, the donor selectors, and
  \code{plateauedEnsembleLevel} are \stub{stubs that return
  false/no-op}. Consequently, no chain is currently cloned mid-cycle to
  accelerate burn-in or to replace a mechanically stuck chain; the
  ensemble relies solely on seeding and adaptation. The \param{cull*},
  \param{stuck*}, and \param{minAcceptanceRate} settings are therefore
  inert at present.
  \item \textbf{ColdStart} (\code{seedColdStart}) currently performs plain
  range-uniform prior draws per chain. The designed importance-resampled
  initialization ($K=\param{coldStartMultiplier}\cdot N_c$ candidates,
  resampled proportional to $\exp(\log\text{-prior}+J)$, with an
  \param{explorerFraction} of raw-prior explorer chains) is \stub{not yet
  implemented}.
  \item \textbf{RatioReseed} (\code{seedRatioWeighted}) and the drift
  detection that would select it are \stub{not implemented}; the caller
  always passes \code{driftDetected = false}, so this mode is never chosen
  and \code{WarmStart}/\code{ProvisionalResume} cover the stationary and
  migrating regimes respectively.
  \item \textbf{Provisional dispersion.} As currently configured, the
  runner \emph{does} publish the partial-pool percentiles for provisional
  cycles (parameter CI history, posterior histograms, and the realization
  band), tagged \code{converged=false}. These are uncertified estimates,
  useful for monitoring while the ensemble migrates; a certified posterior
  width requires a FULL cycle.
  \item \textbf{Inter-cycle stability (Path 2) cannot fire.}
  \code{m\_pointEstimateHistory} is a member of an object rebuilt every
  cycle and receives one push per cycle, so \code{streamStable()} always
  fails its \code{size() >= stabilityWindow} test and returns false. Quorum
  is therefore the sole certification path regardless of
  \param{mcmc\_stability\_enabled}. \stub{Not yet fixed}; it would need the
  ring persisted in the posterior snapshot, exactly as $\hat\Sigma$ and
  $\kappa$ now are.
\end{itemize}

\paragraph{Certification in practice (measured).} In the 244-cycle drift
run of 2026-07, \param{plateaued\_fraction} was \emph{exactly zero in every
cycle} and no cycle certified. The chain is mechanical rather than
statistical: with the proposal scale frozen (the two defects of
Section~``Proposal-scale adaptation''), acceptance sat near $1\%$ against a
target of $0.15$; the motion guard of Section~\ref{sec:plateau} requires
$\param{minAcceptedFraction}\times\param{plateauWindow} = 3$ accepted moves
in the trailing $60$ sweeps, i.e.\ an acceptance of $\ge 5\%$; so chains
flickered between \code{PLATEAUED} and \code{REVERTED} and were classified
climbing whenever the end-of-cycle snapshot was taken. Samples still pooled
from the brief flickers, which is why the point estimates tracked the
drifting truth accurately regardless. \emph{Lowering
\param{mcmc\_quorum\_fraction} cannot address this} --- the input to the
quorum test is already zero; the acceptance rate is the lever.

\paragraph{Source of truth.} The authoritative definitions are
\code{DTStreamingMCMC.h} / \code{DTStreamingMCMC.cpp} (sampler, criteria,
products) and \code{DTAssimilation.cpp} (cycle orchestration, cadence,
output paths). This document should be re-checked against those files when
the stubbed features are implemented.

\end{document}
